% from ega4-17-fr.tex
\oldpage[IV-4]{79}

\subsection{Critères de lissité et non-ramification par fibres}
\EGAReaderTarget{ega4.c04.subsec.17-8}{subsection:IV.17.8}
\label{IV.17.8}

\begin{proposition}[17.8.1]
\label{IV.17.8.1}
Soient $g:Y\to S$, $h:X\to S$ deux morphismes localement de présentation finie. Pour qu'un $S$-morphisme $f:X\to Y$ soit non ramifié, il faut et il suffit que pour tout $s\in S$, le morphisme $f_s:h^{-1}(s)\to g^{-1}(s)$ déduit de $f$ par le changement de base $\Spec(k(s))\to S$, soit non ramifié.
\end{proposition}

Cela résulte aussitôt du fait que, pour un morphisme localement de présentation finie, le fait d'être non ramifié est une propriété des \emph{fibres} de ce morphisme (17.4.1, $d$)), et de ce que, pour tout $y\in Y$, on a $f^{-1}(y)=f_s^{-1}(y)$ si $s=g(y)$.

\begin{proposition}[17.8.2]
\label{IV.17.8.2}
Soient $g:Y\to S$, $h:X\to S$ deux morphismes localement de présentation finie; on suppose en outre que $h$ est plat. Pour qu'un $S$-morphisme $f:X\to Y$ soit lisse (resp. étale), il faut et il suffit que, pour tout $s\in S$, le morphisme $f_s:h^{-1}(s)\to g^{-1}(s)$ déduit de $f$ par le changement de base $\Spec(k(s))\to S$ soit lisse (resp. étale). Lorsqu'il en est ainsi, le morphisme $g$ est plat aux points de $f(X)$.
\end{proposition}

On sait en effet (11.3.10) que pour que $f$ soit plat, il faut et il suffit que $f_s$ le soit pour tout $s\in S$, et qu'alors $g$ est plat aux points de $f(X)$. Mais pour un morphisme plat localement de présentation finie, le fait d'être lisse est une propriété des fibres de ce morphisme (17.5.1, $b$)), et pour tout $y\in Y$, on a $f^{-1}(y)=f_s^{-1}(y)$ si $s=g(y)$.

\begin{remark}[17.8.3]
\label{IV.17.8.3}
Les démonstrations précédentes montrent (compte tenu de (11.3.10)) que si les hypothèses sur $g$ et $h$ sont les mêmes que ci-dessus, alors, pour que $f$ soit non ramifié (resp. lisse, resp. étale) \emph{en un point} $x\in X$, il suffit que, si l'on pose $s=h(x)$, $f_s$ soit non ramifié (resp. lisse, resp. étale) au point $x$.
\end{remark}

\subsection{Morphismes étales et immersions ouvertes}
\EGAReaderTarget{ega4.c04.subsec.17-9}{subsection:IV.17.9}
\label{IV.17.9}

\begin{theorem}[17.9.1]
\label{IV.17.9.1}
Soit $f:X\to Y$ un morphisme. Les propriétés suivantes sont équivalentes:
\begin{enumerate}
  \item[a)] $f$ est une immersion ouverte.
  \item[b)] $f$ est un monomorphisme plat localement de présentation finie.
  \item[c)] $f$ est étale et radiciel.
\end{enumerate}
\end{theorem}

Il résulte de (1.4.3, (i)) que a) implique b). La condition b) implique que pour tout $y\in Y$, la fibre $f^{-1}(y)$ est vide ou isomorphe à $\Spec(k(y))$ (8.11.5.1), donc b) entraîne c), en vertu de (17.6.2, $c$)). Reste à voir que c) implique a).

La question étant locale sur $X$ et sur $Y$ (puisque $f$ est injectif), on peut se borner au cas où $Y$ est affine et $f$ de présentation finie. Comme $f$ est plat, c'est un morphisme \emph{ouvert} (2.4.6), donc, en remplaçant $Y$ par $f(X)$, on peut supposer que $f$ est \emph{surjectif}. Pour tout morphisme $Y'\to Y$, $f'=f_{(Y')}:X_{(Y')}\to Y'$ est encore étale, radiciel, surjectif et de présentation finie, donc ouvert, et par suite un homéomorphisme; autrement dit $f$ est un \emph{homéomorphisme universel}, et étant de type fini et séparé (1.8.7.1), $f$ est \emph{propre}. L'hypothèse que $f$ est radiciel et de type fini entraîne alors que $f$ est quasi-fini; donc (8.11.1) $f$ est un morphisme \emph{fini}. Pour prouver que $f$ est un isomorphisme, on peut se borner au cas

